\specialchapter{Nicolas Bourbaki, Elements of Mathematics}

English translations published:

General Topology, Part I

General Topology, Part II

Theory of Sets

Commutative Algebra

Algebra, Part I

Lie Groups and Lie Algebras, Part I

\textbf{General Topology, Part I}

\textbf{CHAPTER I. TOPOLOGICAL STRUCTURES}

§ 1. Open sets, neighbourhoods, closed sets

§ 2. Continuous functions

§ 3. Subspaces, quotient spaces

§ 4. Product of topological spaces

§ 5. Open mappings and closed mappings

§ 6. Filters

§ 7. Limits

§ 8. Hausdorff spaces and regular spaces

§ 9. Compact spaces and locally compact spaces

§ 10. Proper mappings

§ 11. Connectedness

Exercises

\textbf{CHAPTER II. UNIFORM STRUCTURES}

§ 1. Uniform spaces

§ 2. Uniformly continuous functions

§ 3. Complete spaces

§ 4. Relations between uniform spaces and compact spaces

\textbf{CHAPTER III. TOPOLOGICAL GROUPS}

§ 1. Topologies on groups

§ 2. Subgroups, quotient groups, homomorphisms, homogeneous spaces, product groups

§ 3. Uniform structures on groups

§ 4. Groups operating properly on a topological space; compactness in topological groups and spaces with operators

§ 5. Infinite sums in commutative groups

§ 6. Topological groups with operators; topological rings, division rings and fields

\textbf{§ 7. Inverse limits of topological groups and rings}

\textbf{CHAPTER IV. REAL NUMBERS}

§ 1. Definition of real numbers

§ 2. Fundamental topological properties of the real line

§ 3. The field of real numbers

§ 4. The extended real line

§ 5. Real-valued functions

§ 6. Continuous and semi-continuous real-valued functions

§ 7. Infinite sums and products of real numbers

§ 8. Usual expansions of real numbers; the power of R

Exercises

CHAPTER V. ONE-PARAMETER GROUPS § 1. Subgroups and quotient groups of R § 2. Measurement of magnitudes § 3. Topological characterization of the groups R and T § 4. Exponentials and logarithms Exercises

CHAPTER VI. REAL NUMBER SPACES AND PROJECTIVE SPACES § 1. Real number space R \(^{n}\) § 2. Euclidean distance, balls and spheres § 3. Real projective spaces Exercises

CHAPTER VII. THE ADDITIVE GROUPS R \(^{n}\) § 1. Subgroups and quotient groups of R \(^{n}\) § 2. Continuous homomorphisms of R \(^{n}\) and its quotient groups § 3. Infinite sums in the groups R \(^{n}\) Exercises

CHAPTER VIII. COMPLEX NUMBERS § 1. Complex numbers, quaternions § 2. Angular measure, trigonometric functions § 3. Infinite sums and products of complex numbers § 4. Complex number spaces and projective spaces Exercises

CHAPTER IX. USE OF REAL NUMBERS IN GENERAL TOPOLOGY § 1. Generation of a uniformity by a family of pseudometrics; uniformizable spaces § 2. Metric spaces and metrizable spaces § 3. Metrizable groups, valued fields, normed spaces and algebras § 4. Normal spaces § 5. Baire spaces § 6. Polish spaces, Souslin spaces, Borel sets Appendix: Infinite products in normed algebras Exercises

CHAPTER X. FUNCTION SPACES § 1. The uniformity of S-convergence § 2. Equicontinuous sets § 3. Special function spaces § 4. Approximation of continuous real-valued functions Exercises

\textbf{General Topology, Part II}

\textbf{Theory of Sets}

\textbf{CHAPTER I. DESCRIPTION OF FORMAL MATHEMATICS}

§ 1. Terms and relations

§ 2. Theorems

§ 3. Logical theories

§ 4. Quantified theories

§ 5. Equalitarian theories

Appendix. Characterization of terms and relations Exercises

CHAPTER II. THEORY OF SETS

§ 1. Collectivizing relations

§ 2. Ordered pairs

§ 3. Correspondences

§ 4. Union and intersection of a family of sets

§ 5. Product of a family of sets

§ 6. Equivalence relations

Exercises

CHAPTER III. ORDERED SETS, CARDINALS, INTEGERS

§ 1. Order relations. Ordered sets

§ 2. Well-ordered sets

§ 3. Equipotent sets. Cardinals

§ 4. Natural integers. Finite sets

§ 5. Properties of integers

§ 6. Infinite sets

§ 7. Inverse limits and direct limits

Exercises

CHAPTER IV. STRUCTURES

§ 1. Structures and isomorphisms

§ 2. Morphisms and derived structures

§ 3. Universal mappings

Exercises

SUMMARY OF RESULTS

Introduction

§ 1. Elements and subsets of a set

§ 2. Functions

§ 3. Products of sets

§ 4. Union, intersection, product of a family of sets

§ 5. Equivalence relations and quotient sets

§ 6. Ordered sets

§ 7. Powers. Countable sets

§ 8. Scales of sets. Structures

\textbf{Commutative Algebra}

CHAPTER I. FLAT MODULES § 1. Diagrams and exact sequences § 2. Flat modules § 3. Faithfully flat modules § 4. Flat modules and ``Tor'' functors Exercises

CHAPTER II. LOCALIZATION § 1. Prime ideals § 2. Rings and modules of fractions § 3. Local rings. Passage from the local to the global § 4. Spectra of rings and supports of modules § 5. Finitely generated projective modules. Invertible fractional ideals Exercises

CHAPTER III. GRADUATIONS, FILTRATIONS AND TOPOLOGIES § 1. Finitely generated graded algebras § 2. General results on filtered rings and modules § 3. m-adic topologies on Noetherian rings § 4. Lifting in complete rings § 5. Flatness properties of filtered modules Exercises

CHAPTER IV. ASSOCIATED PRIME IDEALS AND PRIMARY DECOMPOSITION § 1. Prime ideals associated with a module § 2. Primary decomposition § 3. Primary decomposition in graded modules Exercises

CHAPTER V. INTEGERS § 1. Notion of an integral element § 2. The lift of prime ideals § 3. Finitely generated algebras over a field Exercises

CHAPTER VI. VALUATIONS § 1. Valuation rings § 2. Places § 3. Valuations § 4. The height of a valuation § 5. The topology defined by a valuation § 6. Absolute values § 7. Approximation theorem § 8. Extensions of a valuation to an algebraic extension § 9. Application: locally compact fields § 10. Extensions of a valuation to a transcendental extension Exercises

\textbf{COMMUTATIVE ALGEBRA}

CHAPTER VII. DIVISORS

§ 1. Krull domains

§ 2. Dedekind domains

§ 3. Factorial domains

§ 4. Modules over integrally closed Noetherian domains Exercises

\textbf{Algebra, Part I}

\textbf{CHAPTER I. ALGEBRAIC STRUCTURES}

§ 1. Laws of composition; associativity; commutativity

§ 2. Identity element; cancellable elements; invertible elements

§ 3. Actions

§ 4. Groups and groups with operators

§ 5. Groups operating on a set

§ 6. Extensions, solvable groups, nilpotent groups

§ 7. Free monoids, free groups

§ 8. Rings

§ 9. Fields

§ 10. Inverse and direct limits

Exercises

\textbf{CHAPTER II. LINEAR ALGEBRA}

§ 1. Modules

§ 2. Modules of linear mappings. Duality

§ 3. Tensor products

§ 4. Relations between tensor products and homomorphism modules

§ 5. Extension of the ring of scalars

§ 6. Inverse and direct limits of modules

§ 7. Vector spaces

§ 8. Restruction of the field of scalars in vector spaces

§ 9. Affine spaces and projective spaces

§ 10. Matrices

§ 11. Graded modules and rings

Appendix. Pseudomodules

Exercises

\textbf{CHAPTER III. TENSOR ALGEBRAS, EXTERIOR ALGEBRAS, SYMMETRIC ALGEBRAS}

§ 1. Algebras

§ 2. Examples of algebras

§ 3. Graded algebras

§ 4. Tensor products of algebras

§ 5. Tensor algebra. Tensors

§ 6. Symmetric algebras

§ 7. Exterior algebras

§ 8. Determinants

§ 9. Norms and traces

§ 10. Derivations

§ 11. Cogebras, products of multilinear forms, inner products and duality Appendix. Alternative algebras. Octonions

Exercises
